
\documentclass[12pt]{article}
\usepackage[utf8]{inputenc}
\usepackage[T1]{fontenc}
\usepackage[spanish]{babel}
\usepackage{amsmath,amssymb,amsthm}
\usepackage{booktabs}
\usepackage{graphicx}
\usepackage{geometry}
\geometry{margin=2.5cm}

\title{Cambio de Medida y Martingalas en Derivados\\
Verificación de Martingala en Precios Descontados}
\author{Tomas Natividad \and Alvaro Garcia}
\date{Mayo 2026}

\begin{document}

\maketitle

\section{Introducción}

En finanzas matemáticas, uno de los resultados fundamentales de la teoría de valuación libre de arbitraje establece que existe una medida neutral al riesgo $Q$, bajo la cual el precio descontado de un activo riesgoso se comporta como una martingala.

La propiedad central es:

\begin{equation}
E^Q[e^{-rT}S_T]=S_0
\end{equation}

mientras que bajo la medida histórica $P$:

\begin{equation}
E^P[e^{-rT}S_T]=S_0e^{(\mu-r)T}
\end{equation}

El objetivo de este proyecto es verificar numéricamente estas propiedades mediante simulación Monte Carlo utilizando datos reales del ETF SPY.

\section{Marco Teórico}

\subsection{Modelo bajo P}

El activo sigue un Movimiento Browniano Geométrico:

\begin{equation}
\frac{dS_t}{S_t}=\mu dt+\sigma dB_t^P
\end{equation}

cuya solución es:

\begin{equation}
S_T=S_0\exp\left[\left(\mu-\frac{\sigma^2}{2}\right)T+\sigma B_T^P\right]
\end{equation}

\subsection{Cambio de medida}

El precio de mercado del riesgo es:

\begin{equation}
\lambda=\frac{\mu-r}{\sigma}
\end{equation}

y la densidad de Radon-Nikodym:

\begin{equation}
Z_T=\exp\left(-\lambda B_T^P-\frac{1}{2}\lambda^2T\right)
\end{equation}

\subsection{Modelo bajo Q}

\begin{equation}
\frac{dS_t}{S_t}=rdt+\sigma dB_t^Q
\end{equation}

\begin{equation}
S_T=S_0\exp\left[\left(r-\frac{\sigma^2}{2}\right)T+\sigma B_T^Q\right]
\end{equation}

Definiendo el precio descontado:

\begin{equation}
\tilde S_t=e^{-rt}S_t
\end{equation}

se obtiene:

\begin{equation}
d\tilde S_t=e^{-rt}\sigma S_t dB_t^Q
\end{equation}

por lo que $\tilde S_t$ es una martingala bajo $Q$.

\section{Datos y Metodología}

Se utilizaron datos del ETF SPY para el periodo 2021--2023.

\begin{table}[h]
\centering
\begin{tabular}{lc}
\toprule
Parámetro & Valor \\
\midrule
$S_0$ & 462.58 \\
$\mu$ & 11.54\% \\
$\sigma$ & 17.61\% \\
$r$ & 4.50\% \\
$\lambda$ & 0.3999 \\
$N$ & 10,000 \\
$T$ & 1 año \\
\bottomrule
\end{tabular}
\caption{Parámetros utilizados en la simulación}
\end{table}

La simulación Monte Carlo utilizó los mismos choques aleatorios bajo ambas medidas, modificando únicamente el drift.

\section{Resultados}

Los resultados principales fueron:

\begin{align}
E^Q[e^{-rT}S_T] &\approx 461.85 \\
S_0 &= 462.58
\end{align}

y

\begin{equation}
E^P[e^{-rT}S_T]\approx495.54
\end{equation}

confirmando que el precio descontado es martingala bajo $Q$ pero no bajo $P$.

Además,

\begin{equation}
E^P[Z_T e^{-rT}S_T]\approx463.16
\end{equation}

lo cual valida el cambio de medida.

\section{Valuación de Opciones}

El valor de una call europea viene dado por:

\begin{equation}
V_0=e^{-rT}E^Q[(S_T-K)^+]
\end{equation}

Para una opción ATM:

\begin{equation}
V_0 \approx 42.03
\end{equation}

resultado consistente con el precio de Black--Scholes.

\section{Conclusiones}

La simulación Monte Carlo confirmó la propiedad de martingala de los precios descontados bajo la medida neutral al riesgo. Asimismo, se verificó la consistencia del Teorema de Girsanov y la equivalencia entre la valuación Monte Carlo bajo $Q$ y el modelo de Black--Scholes.

\section{Referencias}

\begin{enumerate}
\item Galindo, C. (2025).
\item Black, F. y Scholes, M. (1973).
\item Shreve, S. (2004).
\item Oksendal, B. (2003).
\item Bjork, T. (2009).
\end{enumerate}

\end{document}
